\documentclass{beamer}

\newcommand{\week}{Banks 6 of 6}

\title{Bank Failures}
\subtitle{Mishkin Chapter 10}
\author{Econ 357}
\date{\week}

% Reference the shared preamble
\input{preamble}

% Set the footer -- change 
\setbeamertemplate{footline}{
  \leavevmode%
  \vspace{2ex}
  \hbox{%
    % Left box: Econ 457
    \begin{beamercolorbox}[wd=.4\paperwidth,ht=2.5ex,dp=1ex,left]{author in head/foot}%
      \hspace{1em}Econ 357
    \end{beamercolorbox}%
    % Middle box: Week
    \begin{beamercolorbox}[wd=.2\paperwidth,ht=2.5ex,dp=1ex,center]{date in head/foot}%
      \centering\week
    \end{beamercolorbox}%
    % Right box: Slide numbers
    \begin{beamercolorbox}[wd=.4\paperwidth,ht=2.5ex,dp=1ex,center]{date in head/foot}%
      \hfill\insertframenumber{} 
    \end{beamercolorbox}%
  }%
  \vskip0pt%
}

\begin{document}

\frame{\titlepage}

\begin{frame}
    \frametitle{Outline for Banks}
        \begin{enumerate}
            \item Role of Banks
            \item Bank Balance Sheets
            \item Bank Earnings
            \item Bank Management (and Investment Banks)
            \item Bank Failures
            \item \fbox{Bank Regulation}
        \end{enumerate}
\end{frame}


\begin{frame}
    \frametitle{Outline for Today's Lecture}
    
    \begin{enumerate}

        \item Regulation after the Great Depression
        \begin{itemize}
            \item Glass-Steagall (1933)
            \item Securities and Exchange Commission (1934)
        \end{itemize}
        \item Regulation after the Great Recession
        \begin{itemize}
            \item Stress Tests (2009)
            \item Dodd-Frank (2010)
        \end{itemize}
        \item Moral Hazard and Too Big to Fail
    \end{enumerate}

\end{frame}


\begin{frame}
\frametitle{Notes}

Notes:

    \begin{enumerate}
        \item In the interest of time, we are going to skip the discussion in Mishkin about bank consolidation and international banks.   You can skip Mishkin sections 11.3, 11.4, and 11.7.
        \item Mishkin covers the Great Financial Crisis (GFC) of 2008 in Chapter 12.   We will cover the GFC in depth later in the course.   At this point, we are just discussing the regulation that was enacted after the GFC.  Please read Mishkin sections 12.4 and 12.5.
    \end{enumerate}

\end{frame}


\begin{frame}[t]
\frametitle{1. Regulation after the Great Depression}
\framesubtitle{}

\textbf{Glass-Steagall Act of (1933)}

\begin{itemize}
    \item Created the FDIC
    \item Separation of Commercial Banks and Investment Banks
    \item Commercial banks: take deposits and make loans
    \item Investment banks: underwrote and deal in securities
    \item Intended to make the financial system safer by splitting the riskier investment banking activities from the deposit-based commercial banking activities.
    \item For example, The House of Morgan was split into Morgan Stanley (Investment Bank) and JP Morgan (Commercial Bank).
\end{itemize}

\end{frame}

\begin{frame}[t]
\frametitle{1. Bank Runs}
\framesubtitle{FDIC}

    How to prevent bank runs:  
    \vspace{1em}
    
    FDIC = Federal Deposit Insurance Corporation. Created in 1934 (during the Great Depression)
    \vspace{1em}
    
    Insurance:  Currently the FDIC guarantees up to \$250,000 of deposits.
    \vspace{1em}
    
    Now depositors don’t have incentive to “run” because they know they will get their money back.   Ensures that banks stay in the good equilibrium.

\end{frame}

\begin{frame}[t]
\frametitle{1. Regulation after the Great Depression}
\framesubtitle{FDIC}

    \centering
    \includegraphics[width=0.8\textwidth]{Econ357/figures/03_banks_5_bank_failures.png}

    \blfootnote{Source: Nobel Prize, Scientific Background, 2022}

\end{frame}

\begin{frame}[t]
\frametitle{1. Regulation after the Great Depression}
\framesubtitle{}

    \textbf{Post-script}
    The FDIC remains a very important part of our financial regulatory system.  In fact, it has somewhat expanded it's ability to guarantee short term creditors in banks and it played a very large role in providing guarantees during the 2008 financial crisis.
    \vspace{1em}
    
    In contrast, Glass-Steagall was repealed in 1999 by Gramm-Leach-Bliley.   After that, commercial banks and investment banks were allowed to merge.  This contributed to a trend of consolidation in the banking industry, as can be seen in the next few slides.
    \vspace{1em}

\end{frame}

\begin{frame}[t]
\frametitle{1. Regulation after the Great Depression}
\framesubtitle{Bank Consolidation}

    \centering
    \textbf{Number of Commercial Banks in the US, 1934-2020}
    \vspace{1em}
    
    \includegraphics[width=0.8\textwidth]{Econ357/figures/03_banks_6_consolidation.png}
    
\end{frame}


\begin{frame}[t]
\frametitle{1. Regulation after the Great Depression}
\framesubtitle{Bank Concentration in 2025}

    \footnotesize
    \begin{table}
    \centering
    \begin{tblr}{ 
        colspec = { Q[l,wd=5cm] Q[c,wd=2cm] Q[c,wd=2cm]} 
    }
    \toprule
    Bank &  Assets (\$ tr) &  Share of All Assets (\%) \\
    \midrule
    1. J.P. Morgan Chase \& Co. & 2.74 & 15.4 \\
    2. Bank of America Corp. & 2.38 & 13.4 \\
    3. Citigroup Inc. & 1.96 & 11.0 \\
    4. Wells Fargo \& Co. & 1.89 & 10.6 \\
    5. Goldman Sachs Group &  0.93 &  5.2 \\
    6. Morgan Stanley & 0.88 &  4.9 \\
    7. U.S. Bankcorp & 0.48  & 2.7 \\
    8. P N C Financial Services Group & 0.39 & 2.2 \\
    9. T D Group U S & 0.38 & 2.1 \\
    10. Capital One Financial Corp. & 0.37 & 2.1\\
    \midrule
    Total & 12.40 & 69.6\\
    \bottomrule
    \end{tblr}
    \end{table}
    \end{frame}

\begin{frame}[t]
\frametitle{1. Regulation after the Great Depression}
\framesubtitle{}

    \textbf{Discussion question}: Did the end of Glass-Steagall lead to the financial crisis in 2008?   It's a complicated question with no clear answer.
    \vspace{1em}
    
    On the one hand, the investment side of banks took a lot of risk to generate fees.   At some point this endangered depositors, and the government was forced to step in.   That was precisely what Glass-Steagall had sought to prevent.
    \vspace{1em}
    
    On the other hand, the risk taking likely would have happened anyways, and the investment banks may have required bailouts (the federal government took over AIG, which was an insurer).   Moreover, its not clear that the housing bubble, which was the root cause of the crisis, had anything to do with Glass-Steagall.


\end{frame}


\begin{frame}[t]
\frametitle{1. Regulation after the Great Depression}
\framesubtitle{Securities and Exchange Commission}

    During the “roaring” 1920s it was common for investment houses to engage in what today is considered “securities fraud.”   
     \vspace{1em}
    \begin{itemize}
        \item Charles Ponzi,
        \item “Stock pools”  and pump-and-dump
        \item Misrepresentation of business characteristics
        \item Insider trading
    \end{itemize}
    \vspace{1em}
    
    Many of these were exposed to the public during the Pecora hearings of 1932

\end{frame}

\begin{frame}[t]
\frametitle{1. Regulation after the Great Depression}
\framesubtitle{Securities and Exchange Commission}

The Securities and Exchange Commission (SEC) was created in 1934 in order to regulate securities markets, including both the issuers of securities and the trading of securities.   
\vspace{1em}

At the same time, the Securities Act made it illegal to engage in many of the activities that had caused so much outrage in the Pecora hearings.
\vspace{1em}

The mission of the SEC today is: \textit{“protecting investors, maintaining fair, orderly, and efficient markets, and facilitating capital formation.”}


\end{frame}

\begin{frame}[t]
\frametitle{1. Regulation after the Great Depression}
\framesubtitle{Securities and Exchange Commission}

    \vspace{1em}
    \begin{columns}
    
    \begin{column}{0.5\textwidth}
    Joseph Kennedy Jr. was the first Chairman of the SEC
    \vspace{1em}
    Kennedy had made a fortune in Hollywood and through importing whiskey.
    \vspace{1em}
    
    FDR chose Kennedy, in part, because Kennedy knew the “habits and customs” in the securities industry. (i.e. Kennedy knew about securities fraud)
    
    \end{column}
    
    \begin{column}{0.5\textwidth}
    \centering
    \includegraphics[width=0.7\textwidth]{Econ357/figures/03_banks_6_kennedy.png}
    \end{column}
    
    \end{columns}

\end{frame}

\begin{frame}[t]
\frametitle{1. Regulation after the Great Depression}
\framesubtitle{Securities and Exchange Commission}

    Initially the SEC was small and poorly funded.   It sought to make an impression through large, “show trials” intended to bring the public’s attention to its work.
    \vspace{1em}
    
    Other notable legislation around this time:
    \begin{itemize}
        \item Utility Holding Company Act of 1935 – required regular disclosure of financial condition
        \item Investment Company Act of 1940 – regulates mutual funds
    \end{itemize}

\end{frame}

\begin{frame}[t]
\frametitle{1. Regulation after the Great Depression}
\framesubtitle{Securities and Exchange Commission}

Recent notable SEC cases:

    \begin{itemize}
        \item Michael Milken and Ivan Boesky (1980s)
        \item Enron / WorldCom (2002)
        \item Goldman Sachs (Abacus Deal) (2010)
        \item Bernie Madoff (2008)
        \item Elizabeth Holmes (2018)
        \item Sam Bankman-Fried (2022)
        \item Elon Musk (2026)
    \end{itemize}

\end{frame}

\begin{frame}[t]
\frametitle{1. Regulation after the Great Depression}
\framesubtitle{Securities and Exchange Commission}

SEC actions are generally \textbf{civil} lawsuits.   However, the Department of Justice can also bring \textbf{criminal} charges related to violations of securities laws.

\centering
\includegraphics[width=0.6\textwidth]{Econ357/figures/03_banks_6_criminal.png}

\end{frame}

\begin{frame}[t]
\frametitle{2. Regulation after the Great Recession}
\framesubtitle{Stress Tests}

    March 2009:
        \begin{itemize}
            \item Federal Reserve estimated the amount of capital needed by each of the large banks in order to withstand a shock larger than the Great Depression.
            \item Banks were required to raise enough capital in order to meet that need (by issuing equity, and thereby diluting existing equity holders).
            \item The federal government committed to provide capital to banks, if they weren’t able to raise it privately.
        \end{itemize}

\end{frame}

\begin{frame}[t]
\frametitle{2. Regulation after the Great Recession}
\framesubtitle{Stress Tests}

    \begin{quote}
    By demonstrating that there was going to be enough capital in the American financial system for the system to survive losses like you might face in the Great Depression. We helped make the Great Depression less likely.  \textbf{And another way to say it, we helped put enough money in the window to make the response more credible relative to the scale of potential losses,relative to the scale of potential pressures in the system.} So if you look, you can see, as the broader framework and policy was put in place, people could see it. Get a better feel for it. Its magnitude. Its force. 
    \end{quote}
    - Secretary Tim Geithner

\end{frame}

\begin{frame}[t]
\frametitle{2. Regulation after the Great Recession}
\framesubtitle{Stress Tests}

    \begin{itemize}
        \item Today the Federal Reserve conducts regular stress tests.
        \item They construct a “severely adverse” scenario: GDP -8\%, unemployment rate 10\%, stocks -30\%, VIX 60, etc.
        \item They estimate the losses on the bank balance sheet due to the economic stress.
        \item Large banks are required to maintain enough capital to remain solvent, even if there were to be a large shock.
    \end{itemize}

\end{frame}

\begin{frame}[t]
\frametitle{2. Regulation after the Great Recession}
\framesubtitle{Dodd-Frank Act of 20210}

    \textbf{Dodd-Frank Act of 2010}
    \begin{itemize}
        \item Consumer Financial Protection Bureau
        \item Derivative Pushout: derivatives should be traded on an exchange and cleared
        \item Systemic Risk Regulation: the Fed can designate nonbanks to be “systemic” and subject to regulation (including insurance companies)
        \item Volcker Rule: limit on banks proprietary trading
        \item Resolution Authority
        \item Restrictions on Federal Reserve lending
    \end{itemize}

\end{frame}

\begin{frame}[t]
\frametitle{2. Regulation after the Great Recession}
\framesubtitle{Dodd-Frank Act of 20210}

\textbf{Discussion Question} Did these reforms end Too Big to Fail?   Should they have?
\vspace{1em}

Answer: ?

\end{frame}

\begin{frame}[t]
\frametitle{3. Moral Hazard and Too Big to Fail}
\framesubtitle{}

    Moral hazard is a standard problem in insurance markets.   Insurers may be concerned that providing insurance may incentivize excessive risk taking.
    \vspace{1em}
    
    What do insurers do about this?
    
    \begin{itemize}
        \item Charge different prices based on riskiness of insured (their ability to do this is limited by law, and also because of the Lemons problem and markets could break down).
        \item Monitor behavior, regularly reassess risks
        \item Generally, insurers don’t ever make payments above what was initially contracted.
    \end{itemize}

\end{frame}

\begin{frame}[t]
\frametitle{3. Moral Hazard and Too Big to Fail}
\framesubtitle{Deposit Insurance and Moral Hazard}

Consequence of FDIC guarantee: moral hazard
\vspace{1em}

    \begin{itemize}
        \item Without a guarantee, depositors would have incentives to monitor bank solvency.   They are making a loan to the bank.  They would demand a higher interest rate for higher risk banks.
    
        \item With the guarantee, depositors don’t have to monitor the bank.  Without monitorig, the bank gets loans a low interest rate.   More importantly, after getting the loan (i.e. deposit) the bank may make risky decisions.
    \end{itemize}
    
\end{frame}


\begin{frame}[t]
\frametitle{3. Moral Hazard and Too Big to Fail}
\framesubtitle{}

    \textbf{“Too Big to Fail”}
    \vspace{1em}
    
    \begin{itemize}
        \item During a crisis, it is in the public’s best interest for the government to stop the bank run.   That may require providing insurance ex post, or even bailing out banks, for example by buying assets above “fire sale” prices.
        \item The government may do this, even if there was no insurance contract in place prior to the crisis.
    \end{itemize}
    \vspace{1em}
    
    The 2008 crisis saw many examples of an expansion of insurance, or what were commonly called  “bailouts”

\end{frame}

\begin{frame}[t]
\frametitle{1. Bank Runs}
\framesubtitle{Deposit Insurance and Moral Hazard}

    Consequence of government guarantees more broadly: Too Big to Fail
    
    \begin{itemize}
        \item If a bank is perceived to be “Too Big to Fail” it may cause its managers to take an excess amount of risk.
        \item They get bailed out if they fail, whereas they make money if the risks pay off.
    \end{itemize}

\end{frame}

\begin{frame}[t]
\frametitle{3. Moral Hazard and Too Big to Fail}
\framesubtitle{}

    Moral Hazard / Too Big to Fail concerns may be especially acute in banking:
    
    \begin{itemize}
        \item Depositors have no incentive to monitor risks, because they are covered by FDIC guarantee. Riskier banks don’t have to pay more for deposit funding.
        \item Returns for bank owners/managers are very asymmetric.  They benefit significantly from the upside, where as the downside may be borne by taxpayers.
    \end{itemize}

\end{frame}

\begin{frame}[t]
\frametitle{3. Moral Hazard and Too Big to Fail}
\framesubtitle{}

    \textbf{FDIC Insurance}
    
    \begin{itemize}
        \item All banks are charged fees for being covered by FDIC insurance.
        \item Bank fees are adjusted based on perceived riskiness of the bank.
        \item The fees cover the costs to the taxpayers when the FDIC takes over a bank.   As a result, FDIC generally doesn’t cost the taxpayers money.
        \item However, these fees have not fully addressed the concerns about Moral Hazard.
    \end{itemize}

\end{frame}

\begin{frame}[t]
\frametitle{3. Moral Hazard and Too Big to Fail}
\framesubtitle{}

    \textbf{2008 Bailouts}
    \begin{itemize}
        \item Many banks and insurers received support from the federal government in 2008.
        \item The costs of the bailout in 2008 ended up being very modest, even zero in some calculations.
        \item However, the perception was very damaging.   A clear example was that even after AIG was taken over by the government, many AIG traders were paid multimillion dollar salaries.
        \item Many believed that Wall Street received a bailout, while Main Street did not
    \end{itemize}

\end{frame}

\begin{frame}[t]
\frametitle{3. Moral Hazard and Too Big to Fail}
\framesubtitle{}

    There was a lot of public outrage at banks and at the government after 2008.  
    \vspace{1em}
    
    Question: Was the government justified in bailing out banks?
    \vspace{1em}
    
    Discussion of Matt Taibbi and Robert Shiller
    \vspace{1em}

\end{frame}

\begin{practiceframe}
\frametitle{4. Practice}
\vspace{-1em}

    Which of the following statements is FALSE?
    \begin{enumerate}[A. ]
        \item A bank's assets are its uses of funds.
        \item A bank issues liabilities to acquire funds.
        \item The bank's assets provide the bank with income.
        \item Bank capital is recorded as an asset on the bank balance sheet.
    \end{enumerate}
    
    Which of the following are reported as liabilities on a bank's balance sheet?
    \begin{enumerate}[A. ]
        \item reserves
        \item checkable deposits
        \item consumer loans
        \item deposits with other banks
    \end{enumerate}
    Because [BLANK] are less liquid for the depositor than [BLANK], they earn higher interest rates.
    \begin{enumerate} [A. ]
        \item money market deposit accounts; time deposits
        \item checkable deposits; savings accounts
        \item savings accounts; checkable deposits
        \item savings accounts; time deposits
    \end{enumerate}
    
\end{practiceframe}

\begin{practiceframe}
\frametitle{4. Practice}
\vspace{-1em}

    Banks earn profits by selling [BLANK] with attractive combinations of liquidity, risk, and return, and using the proceeds to buy [BLANK] with a different set of characteristics.
    \begin{enumerate}[A. ]
        \item loans; deposits
        \item securities; deposits
        \item liabilities; assets
        \item assets; liabilities
    \end{enumerate}

    When a \$10 check written on the First National Bank of Chicago is deposited in an account at Citibank, then
    \begin{enumerate}[A. ]
        \item the liabilities of the First National Bank decrease by \$10.
        \item the reserves of the First National Bank increase by \$10.
        \item the liabilities of Citibank decrease by \$10.
        \item the assets of Citibank decrease by \$10.
    \end{enumerate}

\end{practiceframe}

\begin{practiceframe}
\frametitle{4. Practice}
\vspace{-1em}

    For a given return on assets, the lower is bank capital
    \begin{enumerate}[A. ]
        \item the lower is the return for the owners of the bank.
        \item the higher is the return for the owners of the bank.
        \item the lower is the credit risk for the owners of the bank.
        \item the lower the possibility of bank failure
    \end{enumerate}

    Depositors have a strong incentive to show up first to withdraw their funds during a bank crisis because banks operate on a
    \begin{enumerate}[A. ]
        \item last-in, first-out constraint.
        \item sequential service constraint.
        \item double-coincidence of wants constraint.
        \item everyone-shares-equally constraint.
    \end{enumerate}

    Regulators attempt to reduce the riskiness of banks' asset portfolios by
    \begin{enumerate}[A. ]
        \item limiting the amount of loans in particular categories or to individual borrowers.
        \item encouraging banks to hold risky assets such as common stocks.
        \item establishing a minimum interest rate floor that banks can earn on certain assets.
        \item requiring collateral for all loans.
    \end{enumerate}
    



\end{practiceframe}

\end{document}